% Template for ICASSP-2026 paper; to be used with:
%          spconf.sty  - ICASSP/ICIP LaTeX style file, and
%          IEEEbib.bst - IEEE bibliography style file.
% --------------------------------------------------------------------------
\documentclass{article}
\usepackage{spconf,amsmath,graphicx,hyperref}
\usepackage{xcolor}
\usepackage{xspace}
\usepackage{bm}

\newcommand{\xinyue}[1]{\textcolor{purple}{[#1]}}
\newcommand{\ziyu}[1]{\textcolor{blue}{[#1]}}

% ---- system name: change here, changes everywhere ------------------
\newcommand{\sys}{\textsc{Duet}\xspace}

\title{Duet: Coordinating Pretrained Language Models for Joint Generation of Melody and Harmony}

% Three authors, two affiliations
\name{Xinyue Li$^{1}$ \qquad
      Gus Xia$^{1}$ \qquad
      Ziyu Wang$^{1,2}$}

\address{
    $^{1}$Computing and Mathematical Sciences Division, MBZUAI, Abu Dhabi, UAE \\
    $^{2}$Courant Institute, New York University, New York, United States \\
    \texttt{xinyue.li@mbzuai.ac.ae, gus.xia@mbzuai.ac.ae, ziyu.wang@nyu.edu}
}

\begin{document}
\ninept
%
\maketitle
%
\begin{abstract}
Generating two streams together requires each to remain coherent
while accounting for the other at every step. This calls for both
communication and coordination. We introduce \sys, a framework that
enables two pretrained autoregressive models to generate time-aligned
outputs jointly, with each stream informed by the other at every step.
We evaluate \sys on melody--chord generation, where the streams share
a symbolic representation but serve distinct musical roles. This
enables direct comparison with a single model that generates both
under shared weights. Initialized from the same pretrained music
backbone, \sys outperforms the backbone trained from scratch and
achieves subjective ratings comparable to its finetuned single-stream
counterpart, while remaining competitive with external systems. These
findings suggest that coordinated generation can be learned by
composing pretrained models, extending their capabilities beyond
independent generation.
\end{abstract}
%
\begin{keywords}
Multi-stream joint generation, symbolic music generation, model alignment
\end{keywords}
%

\section{Introduction}
\label{sec:intro}

Intelligent behavior often depends on information streams that evolve
together: speech combines text and audio, music coordinates melody
with chords, and embodied behavior links perception with action.
Multimodal models learn relationships between these streams to
interpret their content, for example by answering questions about
images or music, or to generate one stream from another, as in
text-to-speech synthesis. We focus on \emph{coordination}: generating
multiple streams that influence one another and remain consistent as
they unfold. Here, each stream contributes to a shared outcome, rather
than serving primarily as a condition for generating the other.

We study the coordination of independently pretrained language models
by combining two autoregressive models into a single joint generator.
The aim is to let both streams use information from one another while
retaining their pretrained capabilities. This formulation is
particularly relevant when single-stream corpora are abundant but
paired data are scarce: pretraining supplies the individual
capabilities, while paired data support learning their coordination.
More broadly, it provides a setting for studying the composition of
specialized models into a coherent generative system.

We propose \sys, a framework that integrates two pretrained
autoregressive Transformers into a single model for frame-by-frame
co-generation. It combines stream-specific attention, a shared
mixture-of-experts layer, and draft-and-denoise decoding to exchange
information between streams and coordinate their outputs at each step.

We evaluate \sys on melody--chord generation, where the melody and
chord streams are interdependent and share a symbolic representation.
This provides a controlled testbed for comparison with a single model
that generates both streams under shared weights. Both streams in
\sys are initialized from the same pretrained music backbone and
adapted on paired melody--chord data. \sys outperforms the backbone
trained from scratch and achieves subjective ratings comparable to
its finetuned single-stream counterpart, while remaining competitive
with external systems. These results support adapting pretrained
models for coordinated generation with limited paired data.

\section{Related Work}
\label{sec:related}

\textbf{Cross-modal alignment.} Shared embeddings and cross-attention
align modalities for contextual tasks~\cite{clip,flamingo}; conditional
models fix a direction between aligned streams~\cite{yeh2021,sun2021,whole_song}.
Models that emit both streams settle the coupling in advance: unified
models share all weights and learn it through large-scale joint
pretraining~\cite{chameleon,transfusion,bagel}, Mixture-of-Transformers
and MoMa give each modality its own parameters under one global
attention~\cite{mot,moma}, lock-step models such as Moshi~\cite{moshi}
fix an order within each step, and joint diffusion or shared-latent
coupling~\cite{mmdiffusion,codi,seeing_hearing} acts on whole clips.

\textbf{Model merging and mixture of experts.} Sparse upcycling
initializes every expert from one pretrained feed-forward
layer~\cite{upcycling}, and Branch-Train-MiX merges separately trained
experts under a learned router~\cite{btx}; both serve one output stream.
\sys upcycles its pool the same way but learns stream-specific routers
over experts no stream owns, with per-stream attention, for two streams
that must agree at every step.

\section{Methodology}
\label{sec:method}

\begin{figure*}[t]
    \centering
    \includegraphics[width=0.94\linewidth]{figures/arch.pdf}
    \caption{\sys. (a)~Each block attends within a stream (Self Attn)
    and across streams (Cross Attn) with per-stream projections, a
    scalar gate admitting the cross readout. (b)~A router per stream
    selects two of four shared experts, copies of the pretrained
    feed-forward layer. (c)~Joint draft-and-denoise decoding: the
    leader's draft $x^*_t$ is sampled and written in, one denoise
    forward emits $x_t$ and $y_t$ together, and the roles swap at the
    next frame.}
    \label{fig:model}
\end{figure*}

Given autoregressive models $\mathrm{LM}_x$ and $\mathrm{LM}_y$
pretrained independently on streams $x$ and $y$, \sys integrates
them into a joint generator that learns their temporal and semantic
coordination (Fig.~\ref{fig:model}). Stream-specific attention and
routing over shared experts balance specialization with cross-stream
interaction, while draft-and-denoise decoding coordinates each
output pair without fixing either stream as first.

\subsection{Problem Formulation}
\label{sec:formulation}

Let $x=(x_1,\ldots,x_T)$ and $y=(y_1,\ldots,y_T)$ be two co-occurring
streams on a common frame grid, with $z_t=(x_t,y_t)$ denoting the
time-aligned token pair at frame $t$. Aligned generation of $x$ and $y$
seeks to model
\begin{equation}
    p(z_{1:T})
    =
    \prod_{t=1}^{T} p(x_t,y_t \mid z_{<t}),
    \label{eq:joint_ar}
\end{equation}
where both components depend on the shared history $z_{<t}$. A standard
autoregressive model further factorizes each frame as
$p(x_t \mid z_{<t})\,p(y_t \mid z_{<t},x_t)$, or in the reverse order,
making one stream condition on the other's already-generated token.

The problem is therefore how to predict two time-aligned tokens jointly
without fixing either stream as first. \sys interleaves the two
streams: $x_t$ and $y_t$ share rotary position $t$ and predict
$x_{t+1}$ and $y_{t+1}$ respectively, so the next frame is predicted at
two positions of one pass, from one history, rather than one component
after the other. Lastly, the within-frame joint distribution is
realized through a \emph{draft-and-denoise} step,
\begin{equation}
p_\theta(z_t \mid z_{<t})
=
\sum_{d_t}
\underbrace{p_\theta(d_t \mid z_{<t})}_{\text{draft}}\,
\underbrace{p_\theta(x_t \mid z_{<t},d_t)\,p_\theta(y_t \mid z_{<t},d_t)}_{\text{denoise}},
\label{eq:draft_denoise}
\end{equation}
where the draft $d_t$ is the leading stream's own next-frame prediction,
$x^*_t$ or $y^*_t$, the leader alternating across frames, and the two
denoise distributions are read out from one forward pass.

\subsection{Transformer Architecture for Joint Generation}
\label{sec:Duet}

In this section, we discuss important architectural designs in \sys.
Each layer keeps one copy of the pretrained block per stream and aligns
the two in both sub-layers: attention lets the streams read each other
(Fig.~\ref{fig:model}a) and the feed-forward sub-layer routes both over
a shared expert pool (Fig.~\ref{fig:model}b). The attention projections
and the experts are initialized from the pretrained block and
finetuned; the gates and routers are the new parameters.

\textbf{Stream-specific attention}\quad Let $\bm{h}^{(l)}_x, \bm{h}^{(l)}_y$ be
the two streams' representations at layer $l$, and $(W^Q_s, W^K_s,
W^V_s, W^O_s)$ the projections inherited from $\mathrm{LM}_s$. The
information transferred from source stream $a$ to target stream $b$,
$a,b\in\{x,y\}$, is
\begin{equation}
    A_{ab}
    =
    \operatorname{softmax}
    \left(
        \frac{(\bm{h}_b W^Q_b)(\bm{h}_a W^K_a)^\top}{\sqrt{d_k}}
        + M_{ab}
    \right)
    \bm{h}_a W^V_a,
    \label{eq:Duet_attn}
\end{equation}
with queries from the target and keys and values from the source, $d_k$
the head dimension, and $M_{ab}$ a causal mask over frames. $A_{xx}$
and $A_{yy}$ are the self-attention of $\mathrm{LM}_x$ and
$\mathrm{LM}_y$. The cross-stream term reads the partner over earlier
frames and over the current frame, and is combined with the
self-attention through a trainable scalar gate $g_{ab}\in(0,1)$
computed from $\bm{h}_b$, one per direction ($g_{xy}$ and $g_{yx}$);
the current-frame pass is causal because under the shift neither
position holds what it predicts. Then
\begin{equation}
    \bm{o}_b = \bigl(A_{bb} + g_{ab}\,A_{ab}\bigr)\,W^O_b,
    \qquad
    \tilde{\bm{h}}_b = \mathrm{LN}(\bm{h}_b + \bm{o}_b).
    \label{eq:gates}
\end{equation}

\textbf{Shared expert routing in Mixture-of-Experts}\quad
After attention, the feed-forward sub-layer routes both streams
over a shared pool of $E$ experts
$\mathrm{MLP}_1,\ldots,\mathrm{MLP}_E$
(Fig.~\ref{fig:model}b). Each stream scores the pool with its own
router $\boldsymbol{\pi}_s$ and mixes its top-$k$ experts,
\begin{align}
    \boldsymbol{\pi}_s
    &= \operatorname{softmax}(\tilde{\bm{h}}_s G_s),
    \notag \\
    \bm{u}_s
    &= \sum_{i\in\mathcal{T}_k(\boldsymbol{\pi}_s)}
       \pi_{s,i}\,\mathrm{MLP}_i(\tilde{\bm{h}}_s),
    \notag \\
    \bm{h}^{(l+1)}_s
    &= \mathrm{LN}(\tilde{\bm{h}}_s + \bm{u}_s),
    \qquad s\in\{x,y\},
    \label{eq:moe}
\end{align}
where $\mathcal{T}_k$ selects the $k$ largest routing weights, with
the balancing term of~\cite{fedus2022switch}. No expert is assigned
to a stream, so per-stream specialists and shared experts emerge
rather than being imposed. We use $E=4$, $k=2$.

In our experiments $\mathrm{LM}_x$ and $\mathrm{LM}_y$ are the same
pretrained symbolic-music Transformer, so at each layer every expert
starts as a copy of that model's feed-forward network and the four
specialize only through training. When $\mathrm{LM}_x$ and
$\mathrm{LM}_y$ are distinct models with compatible layer structures,
the pool can instead hold the feed-forward networks of both; either
way every expert stays reachable from both routers.

\subsection{Joint Draft-and-Denoise Generation}
\label{sec:generation}

The architecture of Sec.~\ref{sec:Duet} shares the history but
predicts the two components of a frame without seeing each other's
current output. Draft-and-denoise (Eq.~\eqref{eq:draft_denoise})
couples them through a shared draft, refining the frame with the same
network in the spirit of Transfusion's pause-and-denoise
step~\cite{transfusion}, without a separate denoiser.

Two \emph{draft tokens}, one per stream, are appended after the
history at rotary position $t$ (Fig.~\ref{fig:model}c). Each carries a
draft through its stream's frame encoder, or a learned mask embedding
when empty, attends to both histories and to each other, and is never
attended to by history tokens. In training, one frame per window is
chosen at random, its draft tokens filled with model drafts or left
masked, and both are supervised against the true frame alongside the
autoregressive loss. At inference, with $x$ leading at frame $t$, a
first pass predicts both components from the history; only the
leader's is kept and written into its draft token,
$d_t = x^*_t \sim p_\theta(\cdot \mid z_{<t})$, while the follower's
token stays masked. One \emph{denoise} pass then reads out both
components at the draft-token positions, the follower conditioned on
$d_t$ through same-frame attention and the leader re-predicted in place
of its draft, $x_t, y_t \sim p_\theta(\cdot \mid z_{<t}, d_t)$. The
pair is committed, $z_{<t+1} = [z_{<t}, z_t]$, and the leader
alternates at the next frame, so generation stays autoregressive
across frames while both components of a frame come from one pass.

\section{Experimental Setup}
\label{sec:setup_exp}

\subsection{Pretrained backbone}
\label{sec:backbone_training}

Our backbone is the foundation model of~\cite{yinyang}, a
RoFormer~\cite{roformer} with a local encoder that embeds each
sixteenth-note frame, a 12-layer causal decoder of width $d=768$ over
frames, and a local decoder that emits the next frame's notes. It was
pretrained on 357{,}279 files of the Los Angeles MIDI
dataset~\cite{lamd} with all tracks merged into one frame sequence, a
note's track kept only in its program token; we therefore call it
single-stream. \sys copies it into both streams.

\subsection{Training data and training}
\label{sec:data}

We train on POP909~\cite{pop909}, 909 pop songs with a lead melody and
chord annotations; the melody is stream $x$ and the chords, rendered as
a bass note plus up to four upper tones, stream $y$. Every tenth song
is held out (92 after dropping three whose prompt window has no
chords), leaving about 800 for training, far less than the pretraining
corpus, which is why we adapt rather than train from scratch.

\sys keeps the backbone's 12 layers and width, with $E=4$ experts,
$k=2$ active per position and expert width $d_f=4d$. The experts are
initialized as copies of the pretrained feed-forward layer and the
gates at bias $-10$, so training starts from the backbone run
independently on each stream. We train on 384-frame windows with AdamW
and a one-cycle schedule for 15{,}000 steps at a global batch of 128
on 8 GPUs.

\subsection{Baselines}
\label{sec:baselines}

\textbf{Same-backbone} baselines run the backbone as a single stream
over the merged melody and chord sequence, distinguished only by
program tokens, either finetuned on POP909 or trained on it from
scratch; they share \sys's backbone, tokenization and data and isolate
the two-stream construction. \textbf{External} baselines are
Whole-Song Gen~\cite{whole_song}, a hierarchical diffusion model over
form, counterpoint and lead sheet, and the Anticipatory Music
Transformer (AMT)~\cite{anticipatory}, an event-based model that
co-generates both streams but was trained on other data, so it is an
out-of-domain reference rather than a matched opponent.

\section{Evaluation}
\label{sec:eval}

\subsection{Objective evaluation}
\label{sec:objective}

Every system continues the same five-bar prompts of the 92 held-out
POP909 songs to the training length, three samples per song. We
report three groups of metrics: general music quality, prompt
consistency and melody--chord consistency.

\begin{table*}[t]
\centering
\caption{General quality: corpus-pooled divergence from the ground-truth continuation ($\downarrow$, 0 is the reference) over 92 held-out songs; each sample index pooled over all songs against the pooled reference, mean$\pm$std over the three samples. Bold: best per column. $^{*}$: reported but not ranked.}
\label{tab:e1_quality}
\vspace{4pt}
\small
\begin{tabular}{llccccc}
\hline
Group & System & Harm.\ rhythm & Onset grid, mel. & Onset grid, chd. & Duration, mel. & Duration, chd. \\
\hline
\textbf{Ours} & Duet & \textbf{0.0031}$\pm$0.0013 & \textbf{0.0009}$\pm$0.0002 & 0.0018$\pm$0.0009 & 0.0393$\pm$0.0019 & \textbf{0.0032}$\pm$0.0009 \\
 & \hspace{1em}{\footnotesize w/o MoE (dense)} & 0.0042$\pm$0.0004 & 0.0147$\pm$0.0004 & 0.0020$\pm$0.0008 & 0.0366$\pm$0.0009 & 0.0048$\pm$0.0017 \\
 & \hspace{1em}{\footnotesize w/o draft tokens} & 0.0065$\pm$0.0028 & 0.0143$\pm$0.0012 & \textbf{0.0015}$\pm$0.0010 & 0.0380$\pm$0.0019 & 0.0054$\pm$0.0033 \\
\hline
\textbf{Internal} & Single-stream (scratch) & 0.0152$\pm$0.0041 & 0.0016$\pm$0.0002 & 0.0035$\pm$0.0019 & 0.0369$\pm$0.0015 & 0.0112$\pm$0.0042 \\
 & Single-stream (finetuned)$^{*}$ & 0.0031$\pm$0.0006 & 0.0015$\pm$0.0003 & 0.0004$\pm$0.0002 & 0.0353$\pm$0.0012 & 0.0009$\pm$0.0002 \\
\hline
\textbf{External} & Whole-Song Gen & 0.0163$\pm$0.0077 & 0.0039$\pm$0.0003 & 0.0078$\pm$0.0035 & \textbf{0.0070}$\pm$0.0023 & 0.0124$\pm$0.0065 \\
 & AMT & 0.0726$\pm$0.0035 & 0.0025$\pm$0.0011 & 0.0473$\pm$0.0087 & 0.0235$\pm$0.0037 & 0.0814$\pm$0.0082 \\
\hline
\end{tabular}
\end{table*}

\textbf{General quality.}
Table~\ref{tab:e1_quality} reports per-stream Jensen--Shannon
divergence (JSD) between generated and reference distributions of
chord-change intervals, onset positions, and note durations.
The finetuned single-stream model serves as an unranked reference
and remains at the reference floor. Among ranked systems, \sys
leads on three of five metrics and ties on chord onset positions.
Both ablations yield similar chord statistics but increase melody
onset JSD from $0.0009$ to $0.0147$ and $0.0143$. External systems
deviate further in chord-change intervals and chord durations.

\begin{figure*}[t]
  \centering
  \includegraphics[width=0.9\textwidth]{figures/e1_v5b92_prompt_fit.pdf}
  \caption{Objective evaluation on 92 held-out songs, relative to the
  ground-truth continuation ($0$ = matches the real song). Prompt
  consistency: (a)~pitch-class JSD vs.\ prompt, melody; (b)~onset-pattern
  similarity vs.\ prompt, melody; (c,~d)~the same for chord.
  Melody--chord consistency: (e)~chord-tone ratio; (f)~pitch consonance;
  (g)~melody--chord tonal distance; (h)~coupling. Boxes: quartiles over
  songs, whiskers $1.5$\,IQR; $\diamond$ mean; dashed line: reference.
  Systems left of the vertical rule are ranked: solid star, the best of
  them; hollow star, the best of all where different.}
  \label{fig:e1_objective}
\end{figure*}

\textbf{Prompt consistency.} Fig.~\ref{fig:e1_objective}(a--d) measures,
per stream, tonal drift (pitch-class divergence between prompt and
continuation) and rhythmic placement (cosine between their bar-level
onset profiles), each relative to the ground-truth continuation.
Rhythm does not separate the systems, but tonal drift tracks how much
of the pretrained backbone a system keeps: \sys, its no-draft-token
arm and the finetuned single-stream model, which inherit the pretrained
weights intact, drift least and are statistically level, whereas the
dense arm and the model trained from scratch drift three to four times
farther ($p<0.001$) and Whole-Song Gen nearly ten times.

\textbf{Melody--chord consistency.} Fig.~\ref{fig:e1_objective}(e--h)
scores how well the two streams agree: chord-tone ratio and pitch
consonance~\cite{yeh2021}, melody--chord tonal distance, and coupling,
the chord-tone coverage on the true pairing minus that with the chord
track shifted two bars, which only rewards chords that follow
\emph{this} melody. Among systems trained on the same data \sys is
closest to the reference on all four; removing the draft tokens or
training from scratch roughly doubles its gap and one dense block in
place of the experts quadruples it. The finetuned single-stream model
and the external systems sit nearer still, but they are not the target:
the first is the backbone finetuned on the interleaved stream, the
ceiling of joint training on this data, and the others harmonize one
stream given the other, an easier task than writing both at once; they
mark how far co-generation has to go, and \sys covers most of that
distance with both pretrained models intact.

\subsection{Subjective evaluation}
\label{sec:subjective}

Five held-out songs were selected; for each, one sample from \sys and
from each of the four baselines of Sec.~\ref{sec:baselines}, plus the
ground-truth continuation, gives six clips, thirty in all. Each clip
contains the prompt followed by its continuation, all rendered on piano
at one velocity and tempo. Raters heard the groups, and the clips
within a group, in random order and rated every clip from $1$ to $5$
on \emph{prompt consistency}, \emph{melody--chord consistency},
\emph{structure}, \emph{musicality} and \emph{creativity}. 18 raters
(10 male, 8 female; age 21--55; 12 with some musical training, 4 with
none, 2 amateur musicians) gave 402 ratings, 67 per system, in 67
complete rater--song blocks. Ratings were analyzed with a
repeated-measures ANOVA over systems (Greenhouse--Geisser
corrected~\cite{greenhouse1959methods}), followed by paired $t$-tests
against \sys with Holm correction~\cite{holm1979simple}; error bars
are within-subject 95\% confidence
intervals~\cite{cousineau2005confidence,morey2008confidence}.

\begin{figure}[t]
  \centering
  \includegraphics[width=\columnwidth]{figures/subjective_bars_final.pdf}
  \caption{Subjective evaluation: mean rating per system and axis
    (scale $1$--$5$) with within-subject 95\% confidence intervals over
    the 67 rater--song blocks. $*$: differs from \sys, paired $t$-test,
    Holm-corrected, $p < 0.05$.}
  \label{fig:subjective}
\end{figure}

The effect of system is significant on every axis ($p < 10^{-7}$).
Among the generative systems \sys has the highest mean on every axis
but musicality, where the finetuned single-stream model leads ($3.77$
against $3.65$), and overall ($3.81$). It is rated significantly above
the scratch single-stream model and Whole-Song Gen, most on the two
consistency axes and least on musicality and creativity; it does not
differ significantly from AMT or from the finetuned single-stream
model on any axis; and ground truth remains significantly above every
system. The ordering agrees with the prompt and melody--chord
consistency results of Fig.~\ref{fig:e1_objective}.

% \subsection{Expert-routing ablation}
% \label{sec:abl_routing}
%
% We compare four feed-forward designs at equal parameters and activated
% compute, all from the same warm start: a \emph{shared} router scoring
% every position; \emph{stream-specific} routers over one pool (\sys,
% Eq.~\eqref{eq:moe}); a \emph{hard-partitioned} pool, each
% stream restricted to two experts with the balance term computed within
% its pair~\cite{moma}; and a \emph{dense} layer of inner width $8d$,
% matching the activated compute of two $4d$ experts. The two
% stream-specific routers start as one matrix, so at step~0 the model is
% the shared-router model. Routing statistics are read on a fixed
% held-out batch (about 770 positions per stream per layer); a second
% batch reproduces every pattern. [X: per-arm training budget.]
%
% \begin{figure}[t]
%   \centering
%   \begin{minipage}[t]{0.56\columnwidth}
%     \centering
%     \includegraphics[width=\linewidth]{figures/expert_purity.pdf}
%     \\ {\small (a)}
%   \end{minipage}\hfill
%   \begin{minipage}[t]{0.42\columnwidth}
%     \centering
%     \includegraphics[width=\linewidth]{figures/content_probe.pdf}
%     \\ {\small (b)}
%   \end{minipage}
%   \caption{Shared router versus stream-specific routers, same held-out
%     batch. (a)~Expert purity: the share of melody among the positions
%     each expert wins per layer. Poles are per-stream specialists (S,
%     at most $10\%$ or at least $90\%$); the midpoint marks experts
%     serving both streams (I, $40$--$60\%$). (b)~Layers, out of twelve,
%     in which routing inside a stream follows the frame's register
%     beyond a permutation null.}
%   \label{fig:routing}
% \end{figure}
%
% A shared router separates the streams strongly (mean $\ell_1$
% distance $0.59$ of $2$ between the streams' mean routing distributions)
% but mostly by stream identity: with the melody token copied into the
% chord positions, so that only identity differs, $0.41$ remains, and
% swapping the streams' contents moves the expert preferences in none of
% the eleven usable layers. Offering the identity as a per-stream bias on
% the logits changes nothing (bias deltas at most $0.06$); removing it
% from the router's job does, which is what Eq.~\eqref{eq:moe}
% does. Fig.~\ref{fig:routing}(a) shows the result: expert purity, the
% share of melody among the positions an expert wins, rises from nine
% specialist cells of 48 to twelve, no expert being assigned to a stream,
% while some experts keep serving both. Fig.~\ref{fig:routing}(b) shows
% where the freed capacity goes: routing inside a stream follows the
% frame's register above a permutation null in nine chord layers of
% twelve against five for the shared router, while melody stays near the
% ceiling (ten against nine).
%
% \begin{figure}[t]
%   \centering
%   \includegraphics[width=\columnwidth]{figures/load_distribution.pdf}
%   \caption{Distribution of the 48 per-expert top-1 loads (12 layers
%     $\times$ 4 experts) per router design, (a)~teacher-forced on the
%     held-out batch and (b)~free-running, over the routing decisions of
%     the newest positions during generation. Per-layer loads sum to one,
%     so both means are $1/E$ and only the spread differs.}
%   \label{fig:load}
% \end{figure}
%
% Separate routers do not starve experts (Fig.~\ref{fig:load}): none is
% dead in either design, teacher-forced or free-running, and deadliness,
% one minus the mean smallest expert load relative to $1/E$, is $0.25$
% for both under teacher forcing and $0.31$ against the shared router's
% $0.36$ on the model's own output, where the stream-specific routers
% also keep loads closer to balanced ($\sigma = 0.071$ against $0.085$).

\vfill\pagebreak

\bibliographystyle{IEEEbib}
\bibliography{refs}

\end{document}
